\subsection{Native constructors and retained-object transitions}
\label{sec:native-correspondence}

We map the declared record constructors to the fixed Continuity source.
The accompanying audit contains 37 entries across 27 directly cited files,
with exact source ranges and checked Git-blob and SHA-256 identities. It
separates constructor layout, state transition, known-input computation
and the protocol's additional output interfaces. The following correspondences
give the native instances of the objects used in the preceding proofs.

\paragraph{Images, intents and commitments.}
\texttt{durable::seal} encodes and bounds the payload before drawing a
24-byte nonce. The four image tags place that nonce at byte offsets
$[80,104)$, followed by the XChaCha20-Poly1305 ciphertext and authentication
tag. The target-bearing \texttt{PendingWrite} formats have a canonical
tag-dependent prefix, the complete original target image, and a deterministic
HMAC under the journal-derived intent key. These layouts supply the fixed
offset, prefix and intact-target premises of
Lemma~\ref{lem:record-structural-equality}. The wrapping key for a reserved
operation and the intent-authentication key use their respective derivation
domains; the simulator can retain these local storage keys independently of
the selected KEM random tape.

The source uses one image-head domain and five target-bearing intent-digest
domains. Their callers feed the actual returned digests into fences,
checkpoints and subsequent validation. Cancellation intents carry known
descriptors instead of a private target image; selection rows and public
archive descriptors likewise use their actual known-input computations.
This variant distinction supplies the real-versus-designated-private split
used by the commitment table.

\paragraph{Reservation and durable identity.}
Both generation and encapsulation paths persist the reserved operation token
before invoking the corresponding KEM computation. An original retry opens
that token and consumes its retained entropy. The write-intent path first
retains the complete pending wire. Applying it compares the retained original
wire, checks the expected and next heads, installs its original target bytes,
and removes the pending entry. An already installed exact target is handled
as the same operation. A sibling-session write can instead seal a new image
containing unchanged target payload bytes, preserving the target's source
while creating a new physical sealing invocation.

\paragraph{Uncertain commit and recovery.}
The native state distinguishes a caller's outcome from the durable directory.
An unconfirmed commit maps to \texttt{CommitUncertain}; the owner is closed
and subsequently reopened against the retained image and pending target.
An actual committed write therefore persists across an error response.
The explicit storage machine records this outcome and retains the original
operation on resumption. The cited acquisition rule uses the contained
session version's original save event, including when that version is
carried by a later aggregate.

\paragraph{Report dependencies.}
Table~\ref{tab:native-report-families} identifies the three report families.
The resolution key and pending tag are serialized in an epoch's traffic
entry. A later closure or fanout report can authenticate that complete entry.
This gives the causal raw-key edge used by the reverse-order report argument.
The code's exact retry paths retain an original key and tag; a genuinely new
attempt draws a new key. The nonzero check applies to the returned identifier,
with a real zero-ID failure branch. Pending keys use the full 256-bit space.

\begin{table}[t]
\centering\small
\begin{tabularx}{\linewidth}{@{}p{0.17\linewidth}X p{0.25\linewidth}@{}}
\toprule
Family & Immutable transcript and key placement & Ordinary report result \\
\midrule
Session closure & The report MAC covers the session payload with its own
64-byte report/key tail zeroed. The pending key is retained in that tail.
Earlier resolution keys can occur in the preceding traffic entries.
& Session and epoch progress, counts and loss metadata. \\
Fanout abandonment & The MAC covers batch metadata and complete member
payloads. Its own key is held in a separate batch tail, outside that metadata.
Member payloads can include earlier resolution keys.
& Per-member progress and loss metadata. \\
Closed-epoch resolution & The MAC covers the old epoch's counters, message
identifiers, retained digests and skipped indices, excluding its own key.
Its pending key is retained in the traffic entry.
& Includes unconsumed plaintext for the authorized application. \\
\bottomrule
\end{tabularx}
\caption{Source correspondence of the report families. The application's
authorized plaintext output retains its own audience in the security
experiment. Report-key fragments use their actual record content and
origin mapping.}
\label{tab:native-report-families}
\end{table}

The normal retained-object path reopens a registered image and revalidates
its original semantic state. The interface mapping assigns those validation
calls to known state, the protocol interface $F$, or the immutable report
row. Arbitrary modified storage inputs and separately returned projections
receive their corresponding parse, authentication and output operations in
an extended experiment. In particular, the batch-key record and a member's
session payload have different contents: their acquisition events are
determined separately. These typed boundaries preserve the existing target
predicate while making each native call's role in the reduction explicit.
